% Created 2026-02-24 Tue 12:20
% Intended LaTeX compiler: pdflatex
\documentclass[11pt]{article}
\usepackage[utf8]{inputenc}
\usepackage[T1]{fontenc}
\usepackage{graphicx}
\usepackage{longtable}
\usepackage{wrapfig}
\usepackage{rotating}
\usepackage[normalem]{ulem}
\usepackage{amsmath}
\usepackage{amssymb}
\usepackage{capt-of}
\usepackage{hyperref}
\author{Prologos Project}
\date{2026-02-24}
\title{Towards a General Logic Engine on Propagators\\\medskip
\large Lattice-Theoretic Foundations and Architectural Recommendations for Prologos}
\hypersetup{
 pdfauthor={Prologos Project},
 pdftitle={Towards a General Logic Engine on Propagators},
 pdfkeywords={},
 pdfsubject={},
 pdfcreator={Emacs 30.2 (Org mode 9.7.39)}, 
 pdflang={English}}
\begin{document}

\maketitle
\tableofcontents

\section{Executive Summary}
\label{sec:org7243b99}

This document synthesizes two extensive research surveys --- the Lattice-Based
Propagator Networks survey and the Propagator Networks research report --- into
a concrete architectural recommendation for Prologos's logic engine. The
central thesis:

\begin{quote}
Prologos should build its logic engine as a \textbf{stratified propagator network}
with an ATMS worldview layer, using lattice-compatible data structures (LVars)
as the monotonic data plane and assumption-based truth maintenance as the
non-monotonic control plane. Tabling provides completeness; stratification
provides negation. The same substrate should serve both compile-time
(type checking, elaboration) and runtime (logic programming, constraint
solving) use cases.
\end{quote}

The two user-proposed mechanisms --- the "Multiverse Mechanism" (choice-point
forking via ATMS worldviews) and the "Pocket Universe Mechanism" (scoped
computation via lattice embeddings) --- are both \textbf{validated} by the research
and directly implementable. The Multiverse Mechanism maps onto ATMS's
simultaneous worldview maintenance, recovering nondeterminism atop a
monotonic substrate. The Pocket Universe Mechanism maps onto Galois
connections between abstract and concrete domains, enabling modular
constraint domains and scoped sub-computations.

\textbf{Key architectural choices:}

\begin{enumerate}
\item \textbf{Propagator cells} with user-defined lattice merge (Radul-Sussman model)
\item \textbf{ATMS layer} for hypothetical reasoning, \texttt{amb}, and dependency-directed backtracking
\item \textbf{Tabling} for termination and completeness (SLG-style, with lattice answer modes)
\item \textbf{Stratified evaluation} for negation, aggregation, and non-monotonic operations
\item \textbf{Galois-connected domain embeddings} for modular constraint solvers
\item \textbf{Phased implementation} starting with compile-time use (elaborator propagators),
then extending to runtime logic programming
\end{enumerate}
\section{Part I: Lattice Foundations}
\label{sec:org1d4a47e}

\subsection{1. Why Lattices?}
\label{sec:org496ff82}

A \textbf{lattice} is an algebraic structure that captures the notion of
\emph{partial information} and \emph{refinement}. In programming language theory,
lattices appear everywhere:

\begin{itemize}
\item \textbf{Type inference}: the substitution lattice (from most general to most specific)
\item \textbf{Abstract interpretation}: concrete/abstract domain pairs connected by Galois connections
\item \textbf{Constraint propagation}: cell values that accumulate knowledge monotonically
\item \textbf{Distributed systems}: CRDTs converging through lattice joins
\end{itemize}

For a logic engine, lattices provide the mathematical guarantee that
computation \textbf{converges} --- that iteratively refining partial answers
eventually reaches a stable fixed point.
\subsection{2. Crash Course: The Lattice Zoo}
\label{sec:org7910986}

\subsubsection{2.1 Partial Orders}
\label{sec:org92c7757}

A \textbf{partial order} \((S, \leq)\) is a set \(S\) with a binary relation \(\leq\) that is:

\begin{itemize}
\item \textbf{Reflexive}: \(a \leq a\)
\item \textbf{Antisymmetric}: \(a \leq b\) and \(b \leq a\) implies \(a = b\)
\item \textbf{Transitive}: \(a \leq b\) and \(b \leq c\) implies \(a \leq c\)
\end{itemize}

The ordering represents "has at least as much information as." In a type
inference context: \texttt{Int} \(\leq\) \texttt{?a} means "knowing the type is \texttt{Int}
is more informative than knowing nothing (metavariable \texttt{?a})."

Not all elements need be comparable --- this is the "partial" in partial
order. Two type constraints might be independent (neither implies the other).
\subsubsection{2.2 Join-Semilattices}
\label{sec:orgd993916}

A \textbf{join-semilattice} \((L, \leq, \bot, \sqcup)\) adds:

\begin{itemize}
\item A \textbf{bottom} element \(\bot\) (least informative --- "nothing")
\item A \textbf{join} operation \(\sqcup\) (least upper bound --- "combine information")
\end{itemize}

The join is:
\begin{itemize}
\item \textbf{Commutative}: \(a \sqcup b = b \sqcup a\)
\item \textbf{Associative}: \((a \sqcup b) \sqcup c = a \sqcup (b \sqcup c)\)
\item \textbf{Idempotent}: \(a \sqcup a = a\)
\end{itemize}

Join is the \emph{merge} operation in propagator networks. When two propagators
contribute information to the same cell, join combines them. Commutativity
means the order doesn't matter; idempotency means duplicate contributions
are harmless.
\subsubsection{2.3 Complete Lattices and Top}
\label{sec:orgc32f835}

A \textbf{complete lattice} also has:
\begin{itemize}
\item A \textbf{top} element \(\top\) (maximally informative --- \emph{contradiction})
\item \textbf{Meet} operation \(\sqcap\) (greatest lower bound)
\end{itemize}

Top represents incompatible information merged together. When a type
checker discovers that a variable must be both \texttt{Int} and \texttt{String}, the
result is \(\top\) --- contradiction. The propagator network must handle this
by backtracking or reporting an error.
\subsubsection{2.4 Fixed Points: Why Lattices Guarantee Convergence}
\label{sec:org02eac93}

\textbf{Knaster-Tarski theorem}: Every monotone function \(f : L \to L\) on a
complete lattice has a least fixed point:

\(\mathrm{lfp}(f) = \bigsqcap \{ x \in L \mid f(x) \leq x \}\)

\textbf{Kleene iteration}: On a dcpo with \(\bot\), the least fixed point of a
Scott-continuous \(f\) is computed by iterating from bottom:

\(\mathrm{lfp}(f) = \bigsqcup_{n \in \mathbb{N}} f^n(\bot)\)

This is exactly what propagator networks do: start with all cells at \(\bot\)
(no information), iterate propagators (apply \(f\)), and converge to the
least fixed point (the most general solution).
\subsubsection{2.5 Finite Height and Termination}
\label{sec:orgb81fb20}

If the lattice has \textbf{finite height} (longest ascending chain is bounded),
then Kleene iteration terminates in at most \(h\) steps, where \(h\) is the
height. For infinite-height lattices (e.g., intervals over reals),
\textbf{widening} operators accelerate convergence:

\(a \nabla b \geq a \sqcup b\)

Widening over-approximates to ensure termination; \textbf{narrowing} then
recovers precision:

\(a \leq a \triangle b \leq a \quad \text{when } b \leq a\)
\subsection{3. Key Lattice Structures for Logic Programming}
\label{sec:orgf00992d}

\subsubsection{3.1 The Substitution Lattice}
\label{sec:org77cd5e3}

Unification operates on a lattice of substitutions:

\begin{itemize}
\item \(\bot\) = identity substitution (no bindings)
\item \(\sigma \leq \tau\) iff \(\tau\) is a specialization of \(\sigma\) (more bindings)
\item \(\sigma \sqcup \tau\) = most general unifier (mgu) of \(\sigma\) and \(\tau\)
\item \(\top\) = failure (no unifier exists)
\end{itemize}

This is the fundamental lattice for Prolog-style logic programming:
every resolution step computes a join on the substitution lattice.
\subsubsection{3.2 The Herbrand Interpretation Lattice}
\label{sec:org512d620}

For Datalog/bottom-up evaluation:

\begin{itemize}
\item Elements are sets of ground atoms (Herbrand interpretations)
\item \(\bot\) = empty set
\item \(I_1 \leq I_2\) iff \(I_1 \subseteq I_2\)
\item \(I_1 \sqcup I_2 = I_1 \cup I_2\)
\item The immediate consequence operator \(T_P\) maps an interpretation to its
one-step derivation
\end{itemize}

The least fixed point of \(T_P\) gives the minimal Herbrand model --- all
provable ground facts.
\subsubsection{3.3 The Interval Lattice}
\label{sec:orgb43961f}

For numeric constraint propagation:

\begin{itemize}
\item Elements are intervals \([a, b]\) over \(\mathbb{R}\)
\item \(\bot\) = \(\mathbb{R}\) (no constraint)
\item \([a_1, b_1] \sqcap [a_2, b_2] = [\max(a_1, a_2), \min(b_1, b_2)]\)
(intersection = more constrained)
\item \(\top\) = \(\emptyset\) (contradictory constraints)
\end{itemize}
\subsubsection{3.4 The Protocol State Lattice (Session Types)}
\label{sec:org2381593}

For Prologos's session type verification:

\begin{itemize}
\item \(\bot\) = unconstrained channel
\item Interior elements = specific protocol states (e.g., "send Int then receive String")
\item \(\top\) = protocol violation
\item Duality checking: two endpoints' protocol states must be lattice-dual
\end{itemize}
\subsubsection{3.5 The Multiplicity Semiring (QTT)}
\label{sec:org90250d6}

For Prologos's quantitative type theory:

\begin{itemize}
\item Elements: \(\{0, 1, \omega\}\) (erased, linear, unrestricted)
\item Addition: resource usage accumulation
\item Multiplication: usage in composed contexts
\item This forms a \emph{semiring}, not a lattice per se, but the ordering \(0 \leq 1 \leq \omega\)
creates a lattice structure useful for propagation
\end{itemize}
\section{Part II: The Architectural Recommendation}
\label{sec:orgd5aa6b5}

\subsection{4. The Three-Layer Architecture}
\label{sec:orgbb9fa46}

We recommend a three-layer architecture for Prologos's logic engine:

\begin{verbatim}
┌─────────────────────────────────────────────────────┐
│  Layer 3: STRATIFICATION                            │
│  Negation, aggregation, non-monotone operations     │
│  Evaluated stratum-by-stratum (lower strata first)  │
└─────────────┬───────────────────────────┬───────────┘
              │                           │
┌─────────────▼───────────────────────────▼───────────┐
│  Layer 2: ATMS / TRUTH MAINTENANCE                  │
│  Hypothetical reasoning (amb, worldviews)           │
│  Nogood recording, dependency-directed backtracking  │
│  Non-monotonic control over monotonic data           │
└─────────────┬───────────────────────────┬───────────┘
              │                           │
┌─────────────▼───────────────────────────▼───────────┐
│  Layer 1: PROPAGATOR NETWORK                        │
│  Cells (lattice values) + Propagators (monotone fns)│
│  Merge = lattice join                               │
│  Deterministic, order-independent, parallelizable   │
└─────────────────────────────────────────────────────┘
\end{verbatim}
\subsubsection{Layer 1: The Propagator Network (Monotonic Data Plane)}
\label{sec:org1455ada}

The foundation. Cells hold partial information (lattice elements);
propagators compute monotone functions between cells. When a cell's
value increases (via merge/join), all downstream propagators are
scheduled for re-execution. This layer guarantees:

\begin{itemize}
\item \textbf{Determinism}: same inputs \(\to\) same fixed point, regardless of scheduling
\item \textbf{Convergence}: finite-height lattices terminate; widening handles infinite ones
\item \textbf{Parallelism}: propagators can execute concurrently (CALM theorem)
\end{itemize}

This is the layer where type inference, constraint propagation, and
dataflow computation live. It is purely monotonic.
\subsubsection{Layer 2: ATMS / Truth Maintenance (Non-Monotonic Control Plane)}
\label{sec:orge4e0eff}

The innovation. The ATMS maintains \emph{all possible worldviews simultaneously}.
Each derived fact is labeled with the set of assumptions that justify it.

\begin{itemize}
\item \textbf{\texttt{amb}} creates fresh hypothetical premises: \(\{h_1, h_2, \ldots\}\)
\item Each branch's facts are \emph{contingent} on its hypothesis
\item \textbf{Nogoods}: when a contradiction is detected, the responsible assumption set
is recorded as a nogood --- permanently invalidating that worldview
\item \textbf{Dependency-directed backtracking}: nogoods identify \emph{which} choice was
wrong, skipping irrelevant alternatives
\end{itemize}

This layer provides nondeterministic search \emph{without backtracking} in the
traditional sense. All branches exist simultaneously in the ATMS; "exploring
a branch" = believing a hypothesis; "backtracking" = recording a nogood.

Key insight from the research: \textbf{the ATMS's set of nogoods is itself a
monotonically growing lattice}. Non-monotonicity in the search space is
reified as monotonicity in the meta-lattice of learned constraints.
\subsubsection{Layer 3: Stratification (Controlled Non-Monotonicity)}
\label{sec:orgcf440ca}

Some operations are genuinely non-monotonic: negation-as-failure,
aggregation (\texttt{count}, \texttt{min}, \texttt{max}), and finalization. These are handled
by \textbf{stratification}:

\begin{enumerate}
\item Decompose the program into strata (strongly connected components of
the dependency graph, with non-monotonic edges crossing strata)
\item Evaluate strata bottom-up: within each stratum, run the propagator
network to its fixed point
\item Between strata, non-monotonic operations observe the \emph{completed} lower
stratum
\end{enumerate}

This is the established technique from Datalog with negation, Flix,
Bloom\textsuperscript{L}, and CHR. The CALM theorem confirms: within a stratum,
computation is coordination-free; stratum boundaries are the minimal
coordination points.
\subsection{5. Validating the Multiverse Mechanism}
\label{sec:orga314ea8}

The user's proposed "Multiverse Mechanism" --- forking lattice state at
choice points, exploring branches independently --- is \emph{directly validated}
by the research. It maps cleanly onto the ATMS worldview model:

\begin{center}
\begin{tabular}{ll}
Multiverse Concept & ATMS Implementation\\
\hline
Choice point & \texttt{amb} creating hypothetical premises \(\{h_i\}\)\\
Fork & Each \(h_i\) labels a worldview\\
Independent evolution & Facts derived under \(h_i\) are contingent\\
Branch contradiction & Nogood set \(\{h_i, \ldots\}\) recorded\\
Branch success & Answer extracted from worldview \(h_i\)\\
All branches at once & ATMS maintains all worldviews simultaneously\\
\end{tabular}
\end{center}
\subsubsection{Why ATMS over Physical Forking}
\label{sec:org192bd1f}

The research reveals two approaches to forking:

\begin{enumerate}
\item \textbf{Physical forking} (MUSE/OR-parallelism): Copy the entire state at each
choice point. Each copy evolves independently.
\begin{itemize}
\item Pro: Simple, naturally parallel
\item Con: O(state\textsubscript{size}) per fork; no shared learning between branches
\end{itemize}

\item \textbf{Virtual forking} (ATMS): Tag all facts with their supporting assumptions.
All worldviews coexist in a single data structure.
\begin{itemize}
\item Pro: O(1) context switching; nogoods shared across worldviews;
dependency-directed backtracking
\item Con: More complex bookkeeping; label management overhead
\end{itemize}
\end{enumerate}

\textbf{Recommendation for Prologos: ATMS (virtual forking).} The reasons:

\begin{itemize}
\item Prologos already has contingent information in its type checker
(metavariables with constraints that may be retracted)
\item Dependency-directed backtracking produces dramatically better error messages
("this constraint fails because of \emph{these} assumptions") vs chronological
backtracking ("something went wrong somewhere")
\item The learning effect: nogoods discovered in one branch prune other branches.
This is exactly CDCL (Conflict-Driven Clause Learning), which
revolutionized SAT solving
\item Memory efficiency: shared structure between worldviews rather than
full copies
\end{itemize}
\subsubsection{The Multiverse in Practice: A Worked Example}
\label{sec:orgdcb8294}

\begin{verbatim}
;; A Prologos logic program (hypothetical syntax)
relation parent : String -> String -> Prop
parent "alice" "bob"
parent "bob" "carol"
parent "bob" "dave"

relation ancestor : String -> String -> Prop
rule ancestor X Y :- parent X Y
rule ancestor X Y :- parent X Z, ancestor Z Y

query ancestor "alice" ?who
;; Expected: ?who = "bob", ?who = "carol", ?who = "dave"
\end{verbatim}

Under ATMS:
\begin{enumerate}
\item Create cell for \texttt{?who}, initially \(\bot\)
\item First rule: \texttt{parent "alice" ?who} unifies, producing \texttt{?who = "bob"}
under assumption \(h_1\)
\item Second rule: \texttt{parent "alice" ?Z}, \texttt{ancestor ?Z ?who}
\begin{itemize}
\item \texttt{?Z = "bob"} under \(h_2\); then \texttt{ancestor "bob" ?who}
\item First rule: \texttt{parent "bob" ?who} → \texttt{?who = "carol"} under \(h_3\)
or \texttt{?who = "dave"} under \(h_4\) (\texttt{amb} on matching clauses)
\end{itemize}
\item All answers coexist: \texttt{\{h\_1: "bob", h\_3: "carol", h\_4: "dave"\}}
\item Collecting answers = iterating over consistent worldviews
\end{enumerate}
\subsection{6. Validating the Pocket Universe Mechanism}
\label{sec:orgfbb804a}

The "Pocket Universe Mechanism" --- isolating a sub-computation in a
smaller lattice, then projecting results back --- maps onto \textbf{Galois
connections} from abstract interpretation:

\begin{center}
\begin{tabular}{ll}
Pocket Universe Concept & Formal Mapping\\
\hline
Outer lattice (full) & Concrete domain \(C\)\\
Inner lattice (pocket) & Abstract domain \(A\)\\
Enter pocket & Abstraction function \(\alpha : C \to A\)\\
Compute in pocket & Fixed point in \(A\) (cheaper/faster)\\
Exit pocket & Concretization \(\gamma : A \to C\)\\
Soundness guarantee & Galois connection: \(\alpha(x) \leq y \iff x \leq \gamma(y)\)\\
\end{tabular}
\end{center}
\subsubsection{Applications in the Logic Engine}
\label{sec:org921bab0}

\begin{enumerate}
\item \textbf{Tabling as pocket universes}: When a tabled predicate is called, a
sub-computation is launched in a "pocket" (the table). Intermediate
results are accumulated in the table's lattice. When complete, the
table's answers are projected back into the calling context via the
consumer mechanism.

\item \textbf{Modular constraint domains}: Each constraint domain (equality,
ordering, arithmetic, set membership) lives in its own lattice. They
are embedded into the global propagator network via Galois connections.
Each domain solver is a "pocket universe" that receives projections of
the global state and returns refined constraints.

\item \textbf{Scoped hypothetical reasoning}: An \texttt{amb} creates a pocket universe
for each alternative. Within each pocket, propagation proceeds
independently. Nogoods discovered in one pocket prune others via
the ATMS label mechanism.

\item \textbf{Abstract interpretation for type checking}: The type checker operates
in an abstract domain (types as lattice elements) connected to the
concrete domain (values) via a Galois connection. This is already
implicit in Prologos's bidirectional type checking; making it explicit
with propagators would improve composability and error reporting.
\end{enumerate}
\subsection{7. Recovering SLD/SLG Semantics}
\label{sec:orgece61ff}

\subsubsection{7.1 SLD Resolution on Propagators}
\label{sec:orgc9dd85a}

SLD resolution (Prolog's execution model) maps onto propagators as follows:

\begin{center}
\begin{tabular}{ll}
SLD Concept & Propagator Implementation\\
\hline
Goal stack & Worklist of pending propagator firings\\
Clause selection & \texttt{amb} over matching clauses\\
Unification & Join on the substitution lattice\\
New subgoals & New propagators added to the network\\
Backtracking & Nogood recording + hypothesis retraction\\
Answer & Fixed point of the substitution cell\\
\end{tabular}
\end{center}

The key difference: SLD's depth-first, left-to-right search becomes the
ATMS's breadth-first, dependency-directed search. This is \emph{more complete}
(no infinite loops from left-recursion) and \emph{more efficient} (nogoods
prune the search space globally).
\subsubsection{7.2 SLG Resolution (Tabling) on Propagators}
\label{sec:org47d5184}

SLG resolution extends SLD with memoization. The propagator mapping:

\begin{center}
\begin{tabular}{ll}
SLG Concept & Propagator Implementation\\
\hline
Table (memo store) & LVar with lattice answer mode\\
Producer & Propagator computing new answers\\
Consumer & Propagator waiting for table answers (threshold read)\\
Completion & Table LVar frozen; no more answers\\
Well-founded semantics & Stratified evaluation of cyclic dependencies\\
\end{tabular}
\end{center}

\textbf{Lattice answer modes} (from XSB Prolog): instead of storing all
individual answers, the table stores the lattice join of all answers.
This subsumes Datalog evaluation and connects directly to Flix's
lattice-valued relations.

\begin{verbatim}
;; SLG-style tabling with lattice answer mode
;; shortest_path(A, B, D) — D is a MinNat lattice value
table shortest_path : String -> String -> MinNat -> Prop
  :answer-mode lattice  ;; join answers via min

rule shortest_path A B 1 :- edge A B
rule shortest_path A C .{D1 + D2}
  :- edge A B, shortest_path B C D2, .{D1 = 1}
\end{verbatim}
\subsubsection{7.3 Cut and Pruning}
\label{sec:orgf4679a7}

Cut maps onto ATMS operations:

\begin{center}
\begin{tabular}{ll}
Cut Variant & Propagator Implementation\\
\hline
\texttt{!} (hard cut) & Retract alternative hypotheses for this \texttt{amb}\\
\texttt{once/1} & Threshold read + freeze on first answer\\
\texttt{if-then-else} & Conditional propagator (threshold on condition)\\
Committed choice & Guard evaluation + monotone commitment\\
\texttt{\textbackslash{}+} (negation-as-failure) & Stratified: evaluate goal in lower stratum, negate\\
\end{tabular}
\end{center}

The ATMS approach is superior to Prolog's chronological backtracking
because cut's \emph{dependency-directed} counterpart (nogood recording)
naturally identifies \emph{which} choices to prune, avoiding the well-known
pitfalls of "red cuts" that silently change program semantics.

\textbf{Recommendation}: Prologos should support \texttt{once} (deterministic commit)
and \texttt{committed-choice} (guard-based), but NOT raw \texttt{!} (Prolog-style cut).
Cut is too low-level and error-prone; the ATMS provides better
alternatives.
\subsection{8. Tabling: The Completeness Mechanism}
\label{sec:org437c11f}

Tabling is \emph{essential} for a practical logic engine. Without it, left-recursive
rules cause infinite loops (in SLD) or redundant computation (in naive
bottom-up evaluation).
\subsubsection{8.1 Design for Prologos}
\label{sec:orgda551bd}

\begin{verbatim}
;; Tabling declaration via spec metadata
spec ancestor : String -> String -> Prop
  :tabled true
  :answer-mode all  ;; or :answer-mode lattice for aggregate answers

;; Implementation is just regular rules
rule ancestor X Y :- parent X Y
rule ancestor X Y :- parent X Z, ancestor Z Y
\end{verbatim}

Under the hood:
\begin{enumerate}
\item First call to \texttt{ancestor} creates a \emph{producer} propagator
\item The producer evaluates rules, pushing answers into a table (LVar)
\item Subsequent calls create \emph{consumer} propagators that read from the table
via threshold reads
\item When the producer reaches a fixed point (no new answers), the table is frozen
\item Consumers receive all answers
\end{enumerate}
\subsubsection{8.2 Lattice Answer Modes}
\label{sec:orgd8e6395}

Following XSB Prolog's design, each tabled predicate can specify an answer
mode:

\begin{itemize}
\item \textbf{\texttt{all}}: Store all distinct answers (standard tabling). The lattice is
the powerset of substitutions.
\item \textbf{\texttt{lattice(f)}}: Aggregate answers via a user-specified lattice join \texttt{f}.
New answers are only "new" if they improve the current aggregate.
This subsumes Datalog with lattices (Flix-style).
\item \textbf{\texttt{first}}: Store only the first answer (\texttt{once} semantics, table frozen
after first answer).
\end{itemize}
\subsubsection{8.3 Tabling as Pocket Universe}
\label{sec:org9bcefca}

Each tabled predicate's table is a "pocket universe" in the Galois
connection sense:

\begin{itemize}
\item \textbf{Abstraction}: project the calling context's substitution onto the
tabled predicate's argument positions
\item \textbf{Computation}: evaluate the tabled predicate's rules within the pocket
\item \textbf{Concretization}: lift the table's answers back into the calling context
\end{itemize}

This ensures that tabled sub-computations are modular and can be
incrementally updated when the calling context provides new information.
\subsection{9. Connecting to Existing Prologos Infrastructure}
\label{sec:org8a52c01}

\subsubsection{9.1 The Elaborator as a Propagator Network (Today)}
\label{sec:org982298b}

Prologos's elaborator \emph{already} implements propagator-like behavior:

\begin{center}
\begin{tabular}{ll}
Current Elaborator & Propagator Equivalent\\
\hline
Metavariables & Cells (lattice of types/terms)\\
Unification constraints & Propagators (join on subst lattice)\\
\texttt{resolve-trait-constraints!} & Propagators resolving dict params\\
\texttt{save/restore-meta-state!} & Speculative execution (worldviews)\\
\texttt{check-unresolved} & Detecting unsolved cells\\
Zonking & Reading final fixed-point state\\
\end{tabular}
\end{center}

The gap: these are \emph{ad hoc} rather than \emph{architected}. The metavar store
is a mutable hash table, not a lattice-structured cell. Trait resolution
uses a worklist, but doesn't record nogoods or do dependency-directed
backtracking. Speculative execution (\texttt{save/restore-meta-state!}) copies
the entire state, rather than using ATMS-style contingent values.

\textbf{Phase 1 recommendation}: Refactor the elaborator's metavar system to
use proper propagator cells with lattice merge semantics. This would
immediately improve:
\begin{itemize}
\item Error messages (dependency tracking shows \emph{why} a constraint failed)
\item Implicit argument inference (more information from bidirectional flow)
\item Speculative type checking (ATMS worldviews instead of full state copy)
\end{itemize}
\subsubsection{9.2 Session Types as Protocol Propagation}
\label{sec:orgc0c751a}

Session type checking already propagates protocol state. Making this
explicit with propagator cells would enable:
\begin{itemize}
\item Multi-party session typing (currently only binary)
\item Dependent session types (protocol depends on values)
\item Better error messages for protocol violations
\end{itemize}
\subsubsection{9.3 QTT Multiplicity Checking as Propagation}
\label{sec:org8c0a660}

Multiplicity inference (Sprint 7) already propagates multiplicity
constraints. With proper propagator cells:
\begin{itemize}
\item Multiplicity inference becomes a standard fixpoint computation
\item Multiplicity errors get dependency tracking
\item The semiring structure (\(0, 1, \omega\)) is naturally a lattice
\end{itemize}
\section{Part III: The Concrete Architecture}
\label{sec:org38611d2}

\subsection{10. Core Data Types}
\label{sec:org96ba619}

\subsubsection{10.1 Lattice Trait}
\label{sec:orgcfd006b}

\begin{verbatim}
;; The fundamental lattice abstraction
trait Lattice {A : Type}
  bot   : A
  join  : A -> A -> A
  leq   : A -> A -> Bool

;; Laws (checked via property testing):
;;   join is commutative:  join a b = join b a
;;   join is associative:  join (join a b) c = join a (join b c)
;;   join is idempotent:   join a a = a
;;   bot is identity:      join bot a = a
;;   leq is consistent:    leq a b = (join a b == b)
\end{verbatim}
\subsubsection{10.2 Propagator Cell}
\label{sec:orgb13f8b1}

\begin{verbatim}
;; A cell holds a lattice value, tracks dependents
deftype Cell (A : Type)
  :where (Lattice A)
  mk-cell : A -> Cell A

;; Core operations
spec cell-read  : {A : Type} where (Lattice A) [Cell A] -> A
spec cell-write : {A : Type} where (Lattice A) [Cell A] -> A -> Unit
  ;; cell-write c v = merge(current, v); schedule dependents if changed

spec cell-watch : {A : Type} where (Lattice A) [Cell A] -> [A -> Unit] -> Unit
  ;; Register a callback (propagator) triggered on cell value change
\end{verbatim}
\subsubsection{10.3 Supported Values (ATMS Integration)}
\label{sec:org468f5dd}

\begin{verbatim}
;; A supported value pairs information with its justification
deftype Supported (A : Type)
  mk-supported : A -> [Set Assumption] -> Supported A

;; A TMS cell holds multiple contingent values
deftype TMSCell (A : Type)
  :where (Lattice A)
  mk-tms-cell : [Set [Supported A]] -> TMSCell A
\end{verbatim}
\subsubsection{10.4 LVar (Lattice Variable)}
\label{sec:org398da1e}

\begin{verbatim}
;; LVar: monotonic lattice variable with threshold reads
deftype LVar (A : Type)
  :where (Lattice A)
  mk-lvar : A -> LVar A

spec lvar-put : {A : Type} where (Lattice A)
  [LVar A] -> A -> Unit
  ;; Inflationary: new_state = join(current, value)

spec lvar-get : {A : Type} where (Lattice A)
  [LVar A] -> [Set A] -> A
  ;; Threshold read: block until state >= some element of threshold set
  ;; Threshold elements must be pairwise incompatible

spec lvar-freeze : {A : Type} where (Lattice A)
  [LVar A] -> A
  ;; Freeze and return exact state (quasi-deterministic)
\end{verbatim}
\subsection{11. The Propagator Network Runtime}
\label{sec:orgf41ca3b}

\subsubsection{11.1 Scheduler}
\label{sec:org6aa6b4c}

\begin{verbatim}
┌──────────────────────────────┐
│       Scheduler              │
│  ┌───────────────────────┐   │
│  │  Worklist (queue)     │   │
│  │  [prop1, prop3, ...]  │   │
│  └───────────────────────┘   │
│                              │
│  Loop:                       │
│    1. Dequeue propagator     │
│    2. Execute (read inputs)  │
│    3. Write outputs (merge)  │
│    4. If cell changed:       │
│       enqueue dependents     │
│    5. If cell = top:         │
│       contradiction handler  │
│    6. Repeat until empty     │
└──────────────────────────────┘
\end{verbatim}

Properties:
\begin{itemize}
\item Jobs are idempotent (safe to re-execute)
\item Scheduling order doesn't affect the final result (lattice determinism)
\item Parallelizable: independent propagators can execute concurrently
\end{itemize}
\subsubsection{11.2 Contradiction Handler}
\label{sec:orgb09b769}

When a cell reaches \(\top\) (contradiction):

\begin{enumerate}
\item Extract the \emph{support set} (assumptions that led to contradiction)
\item Record as a \textbf{nogood} in the ATMS
\item Retract the most recently chosen hypothesis in the nogood set
\item Re-propagate with the retracted hypothesis disbelieved
\end{enumerate}

This is dependency-directed backtracking. It avoids the exponential
blowup of chronological backtracking by targeting the actual cause.
\subsubsection{11.3 The \texttt{amb} Operator}
\label{sec:org0251a8e}

\begin{verbatim}
;; amb creates a choice point with n alternatives
spec amb : {A : Type} [List A] -> A

;; Under the hood:
;; 1. Create fresh hypotheses h1, h2, ..., hn
;; 2. For each hi, create supported value (ai, {hi})
;; 3. Add mutual exclusion: exactly one hi is believed
;; 4. Return the cell containing the contingent value
\end{verbatim}
\subsubsection{11.4 Answer Collection}
\label{sec:org9206a95}

\begin{verbatim}
;; Collect all solutions from a nondeterministic computation
spec solve-all : {A : Type} [Unit -> A] -> [List A]

;; Under the hood:
;; 1. Run the computation with ATMS
;; 2. Iterate over all consistent worldviews at quiescence
;; 3. For each worldview, extract the answer
;; 4. Return list of all answers (duplicates removed)
\end{verbatim}
\subsection{12. Phased Implementation Plan}
\label{sec:org42751be}

\subsubsection{Phase 0: Lattice Trait + Basic Cells (Foundation)}
\label{sec:org485a968}

\textbf{Goal}: Establish the lattice trait and basic propagator cell infrastructure
at the Racket level.

\begin{itemize}
\item Define \texttt{Lattice} trait in Prologos
\item Provide instances: \texttt{FlatLattice} (bottom/value/top), \texttt{SetLattice}
(powerset with union), \texttt{IntervalLattice}
\item Implement \texttt{Cell} as a Racket-level struct with mutable value and
dependent list
\item Implement basic scheduler (worklist, fire-till-quiescence)
\item \textbf{\textasciitilde{}15 AST nodes} (cell-new, cell-read, cell-write, cell-watch, \ldots{})
\item \textbf{\textasciitilde{}40 tests}
\end{itemize}

\textbf{Dependencies}: None (builds on existing trait system)
\textbf{Estimated effort}: Medium
\subsubsection{Phase 1: Elaborator Propagators (Compile-Time Use)}
\label{sec:orgf29a4dd}

\textbf{Goal}: Refactor the metavar system to use propagator cells internally.

\begin{itemize}
\item Replace \texttt{current-meta-store} with propagator cells
\item Unification constraints become propagators between type cells
\item Trait resolution constraints become propagators
\item Add dependency tracking to metavar refinements
\item Improve error messages using dependency information
\item \textbf{No new surface syntax} --- this is internal refactoring
\item \textbf{\textasciitilde{}60 tests} (regression + new dependency-tracking tests)
\end{itemize}

\textbf{Dependencies}: Phase 0
\textbf{Estimated effort}: Large (touches elaborator + typing-core + unify)
\textbf{Key risk}: This is a significant refactoring of the elaborator core
\subsubsection{Phase 2: LVars + Tabling (Runtime Foundation)}
\label{sec:org7fc556d}

\textbf{Goal}: LVar data structures and tabled evaluation.

\begin{itemize}
\item Implement \texttt{LVar} with \texttt{put}, threshold \texttt{get}, and \texttt{freeze}
\item Implement \texttt{LVar-Set}, \texttt{LVar-Map} (grow-only collections)
\item Implement tabling framework: producer/consumer pattern, table LVars,
completion detection
\item Add \texttt{:tabled} spec metadata handling
\item \textbf{\textasciitilde{}20 AST nodes} (lvar-new, lvar-put, lvar-get, lvar-freeze, table-lookup, \ldots{})
\item \textbf{\textasciitilde{}50 tests}
\end{itemize}

\textbf{Dependencies}: Phase 0, Phase 2d of Core DS roadmap (transient builders)
\textbf{Estimated effort}: Large
\subsubsection{Phase 3: ATMS Layer (Nondeterministic Search)}
\label{sec:orgc29260c}

\textbf{Goal}: Assumption-based truth maintenance for hypothetical reasoning.

\begin{itemize}
\item Implement \texttt{Assumption}, \texttt{Supported}, \texttt{TMSCell} types
\item Implement nogood recording and management
\item Implement \texttt{amb} operator with hypothesis creation
\item Implement contradiction handler with dependency-directed backtracking
\item Implement \texttt{solve-all} answer collection
\item \textbf{\textasciitilde{}15 AST nodes} (amb, assume, retract, nogood, \ldots{})
\item \textbf{\textasciitilde{}60 tests}
\end{itemize}

\textbf{Dependencies}: Phase 0, Phase 2
\textbf{Estimated effort}: Large (most complex phase)
\subsubsection{Phase 4: Stratified Evaluation (Negation + Aggregation)}
\label{sec:orgb5ed10a}

\textbf{Goal}: Support negation-as-failure and aggregation via stratification.

\begin{itemize}
\item Implement SCC decomposition of rule dependency graphs
\item Implement stratum-by-stratum evaluation
\item Support \texttt{not} (negation-as-failure) between strata
\item Support lattice aggregation (\texttt{count}, \texttt{min}, \texttt{max}, \texttt{sum}) between strata
\item \textbf{\textasciitilde{}10 AST nodes}
\item \textbf{\textasciitilde{}40 tests}
\end{itemize}

\textbf{Dependencies}: Phase 3
\textbf{Estimated effort}: Medium
\subsubsection{Phase 5: Surface Syntax + Integration (User-Facing Logic Programming)}
\label{sec:org2032bc7}

\textbf{Goal}: Prologos surface syntax for logic programming.

\begin{itemize}
\item \texttt{relation} declarations
\item \texttt{rule} and \texttt{query} syntax
\item Pattern matching integration with unification
\item \texttt{defn} bodies that can contain logic variables
\item Integration with the prelude and module system
\item \textbf{Grammar updates} (grammar.ebnf, grammar.org)
\item \textbf{\textasciitilde{}80 tests}
\end{itemize}

\textbf{Dependencies}: All previous phases
\textbf{Estimated effort}: Large
\subsubsection{Phase 6: Galois Connections + Domain Embeddings (Advanced)}
\label{sec:orgcbc91ce}

\textbf{Goal}: Modular constraint domains connected via Galois connections.

\begin{itemize}
\item \texttt{domain} declarations with abstraction/concretization functions
\item Cross-domain propagation
\item Abstract interpretation framework for program analysis
\item \textbf{\textasciitilde{}30 tests}
\end{itemize}

\textbf{Dependencies}: Phase 5
\textbf{Estimated effort}: Medium-Large
\section{Part IV: Critique and Open Questions}
\label{sec:org84382c5}

\subsection{13. What This Architecture Does Well}
\label{sec:org8119818}

\begin{enumerate}
\item \textbf{Unification of paradigms}: The same propagator substrate serves type
checking, constraint solving, and logic programming. This is not merely
aesthetic --- it means improvements to the propagator infrastructure
benefit all three use cases.

\item \textbf{Sound theoretical basis}: Every component has well-understood
lattice-theoretic foundations. Convergence, determinism, and soundness
are not hoped for but \emph{proven} by the mathematical framework.

\item \textbf{Incremental by construction}: Propagator networks naturally support
incremental computation. When a constraint is added or retracted, only
affected cells re-propagate. This is essential for interactive
development (IDE integration, REPL).

\item \textbf{Parallelizable by the CALM theorem}: Within a stratum, all
propagation is monotonic and therefore coordination-free. This
provides a natural path to parallel type checking and parallel
logic evaluation.

\item \textbf{Better error messages}: Dependency tracking in the ATMS means every
derived fact carries its provenance. Type errors can show exactly
\emph{which} constraints are in tension and \emph{where} they originated.
\end{enumerate}
\subsection{14. What This Architecture Doesn't Solve (Yet)}
\label{sec:org9fa9c18}

\begin{enumerate}
\item \textbf{Performance}: ATMS label management has overhead. For simple
deterministic programs, the ATMS layer is pure cost. We need a
\emph{bypass} for deterministic code paths that avoids the TMS entirely.

\item \textbf{Infinite domains}: Widening operators are needed for lattices of
infinite height (intervals, numeric constraints). Choosing the right
widening strategy is domain-specific and sometimes tricky.

\item \textbf{Efficiency of lattice answer modes}: When the lattice is large or
the join is expensive, tabling with lattice answer modes can be slow.
XSB Prolog's experience shows this requires careful implementation.

\item \textbf{Surface syntax design}: How should logic programming syntax look in
Prologos? The document proposes \texttt{relation} / \texttt{rule} / \texttt{query} but
this needs design iteration. Should it integrate with \texttt{match}? With
\texttt{where} constraints? With mixfix \texttt{.\{...\}}?

\item \textbf{Interaction with QTT}: Logic variables are inherently shared
(multiple references). How does this interact with linear types?
Preliminary answer: logic variables live at multiplicity \(\omega\)
(unrestricted), but the \emph{binding environment} can be linear.
\end{enumerate}
\subsection{15. Open Research Questions}
\label{sec:org473d69f}

\begin{enumerate}
\item \textbf{Can ATMS label management be made zero-cost for deterministic programs?}
The "speculative ATMS" idea: don't create labels until the first \texttt{amb}.
Deterministic programs never pay the ATMS overhead.

\item \textbf{How do Prologos's dependent types interact with tabling?}
Tabling requires comparing calls for "variance" (same query up to
renaming). With dependent types, this comparison involves type equality,
which may itself require propagation. Circularity risk.

\item \textbf{Can propagator-based elaboration subsume the current ad-hoc metavar
system without performance regression?} The current system is fast
because it's specialized. A general propagator system may be slower
without careful optimization.

\item \textbf{What is the right granularity for stratification?}
Per-predicate? Per-SCC? Per-module? Finer granularity allows more
parallelism but increases analysis cost.

\item \textbf{How do pocket universes compose?} Can a tabled predicate call another
tabled predicate? (Yes, but completion detection becomes more complex ---
this is the "scheduling dependency" problem in SLG resolution.)
\end{enumerate}
\section{Summary of Recommendations}
\label{sec:org23d9433}

\begin{center}
\begin{tabular}{rll}
\# & Recommendation & Confidence\\
\hline
1 & Radul-Sussman propagator model as the substrate & High\\
2 & ATMS for hypothetical reasoning (validates Multiverse Mech.) & High\\
3 & Galois connections for domain modularity (validates Pocket U.) & High\\
4 & SLG-style tabling with lattice answer modes & High\\
5 & Stratified evaluation for negation and aggregation & High\\
6 & Start with compile-time use (elaborator), then extend to runtime & High\\
7 & No raw Prolog-style \texttt{!} cut; use \texttt{once} and committed choice & Medium\\
8 & "Speculative ATMS" for zero-cost deterministic bypass & Medium\\
9 & \texttt{Lattice} trait as foundation (user-extensible domains) & High\\
10 & Surface syntax: \texttt{relation} / \texttt{rule} / \texttt{query} keywords & Low (needs design)\\
\end{tabular}
\end{center}

\textbf{Overall assessment}: The research strongly supports the proposed
architecture. The Multiverse and Pocket Universe mechanisms are not
speculative --- they are well-established techniques (ATMS worldviews
and Galois connections respectively) with decades of theoretical
grounding and practical implementation experience. The main risk is
implementation complexity, mitigated by the phased approach starting
with the elaborator refactoring.
\section{References}
\label{sec:org7622514}

Organized by relevance to the recommended architecture:
\subsection{Foundational (Must-Read)}
\label{sec:org7278eb0}
\begin{itemize}
\item Radul \& Sussman, "The Art of the Propagator" (MIT TR, 2009)
\item Radul, \emph{Propagation Networks} (PhD thesis, MIT, 2009)
\item de Kleer, "An Assumption-Based TMS" (\emph{AI Journal}, 1986)
\item Kuper \& Newton, "LVars: Lattice-Based Data Structures for Deterministic Parallelism" (FHPC, 2013)
\item Kuper et al., "Freeze After Writing" (POPL, 2014)
\end{itemize}
\subsection{Tabling and Logic Programming}
\label{sec:orga7795f7}
\begin{itemize}
\item Swift \& Warren, "XSB: Extending Prolog with Tabled Logic Programming" (\emph{TPLP}, 2012)
\item Chen \& Warren, "Tabled Evaluation with Delaying" (\emph{JACM}, 1996)
\item Gupta et al., "Parallel Execution of Prolog Programs: A Survey" (\emph{TOPLAS}, 2001)
\end{itemize}
\subsection{Lattice-Based Languages}
\label{sec:org53c296e}
\begin{itemize}
\item Madsen et al., "From Datalog to Flix" (PLDI, 2016)
\item Arntzenius \& Krishnaswami, "Datafun: A Functional Datalog" (ICFP, 2016)
\item Arntzenius \& Krishnaswami, "Seminaive Evaluation for a Higher-Order Functional Language" (POPL, 2020)
\end{itemize}
\subsection{Theoretical Foundations}
\label{sec:org03cc6b8}
\begin{itemize}
\item Tarski, "A Lattice-Theoretical Fixpoint Theorem" (\emph{Pacific J. Math.}, 1955)
\item Cousot \& Cousot, "Abstract Interpretation: A Unified Lattice Model" (POPL, 1977)
\item Denecker et al., "Approximation Fixpoint Theory" (2012)
\item Hellerstein, "Keeping CALM" (\emph{CACM}, 2020)
\end{itemize}
\subsection{Constraint Systems}
\label{sec:org9844d6a}
\begin{itemize}
\item Fruhwirth, \emph{Constraint Handling Rules} (CUP, 2009)
\item Shapiro et al., "Conflict-Free Replicated Data Types" (INRIA, 2011)
\item Conway et al., "Logic and Lattices for Distributed Programming" (SoCC, 2012)
\end{itemize}
\subsection{Type Systems}
\label{sec:org4f7f2a1}
\begin{itemize}
\item Dunfield \& Krishnaswami, "Bidirectional Typechecking for Higher-Rank Polymorphism" (ICFP, 2014)
\item Pottier \& Remy, "The Essence of ML Type Inference" (2005)
\item Kovacs, \emph{Elaboration Zoo} (2022)
\end{itemize}
\subsection{Implementations}
\label{sec:org5936c25}
\begin{itemize}
\item Holmes (Haskell): \texttt{github.com/i-am-tom/holmes}
\item Propaganda (Clojure): \texttt{github.com/tgk/propaganda}
\item Ascent (Rust): \texttt{github.com/s-arash/ascent}
\item Guanxi (Haskell): \texttt{github.com/ekmett/guanxi}
\end{itemize}
\end{document}
